% !TeX root = closed-systems-from-comparison-completeness.tex
% ============================================================
\chapter{Matter-Sector Realization of the Canonical Scalar Invariant}
\label{ch:matter-realization}
% ============================================================

% ============================================================
\section{Purpose and logical status}
\label{sec:matter-realization-purpose}
% ============================================================

\Cref{ch:numbers-shell-closure} closes the intrinsic numerical classification
of the scalar channel at the level of quartic leading behavior and a
single normalization orbit, while recording same-channel correction
structure only as compatibility-level bookkeeping.
What remains is not further intrinsic classification, but realization role.
For the Chapter~21 mass-coordinate reformulation and the accompanying
conditional exact-ratio package used here as downstream numerical
bookkeeping, see
\cref{sec:intrinsic-spectral-mass-law}.

This chapter identifies that role for the canonical excitation-sector scalar
invariant.
No observer-side measured-mass interpretation is assumed here.
The objective is purely internal: determine which already-established
realization branch canonically carries the Chapter~19 sector invariant.

The argument unfolds in four steps:
\begin{enumerate}[label=\textup{(\roman*)}]
\item type the intrinsic scalar datum and its realization roles;
\item prove branch priority and the intrinsic/observer boundary;
\item isolate the operator-realization gap and its form route;
\item formalize observer-map factorization constraints.
\end{enumerate}

% ============================================================
\section{Structural inputs}
\label{sec:matter-realization-inputs}
% ============================================================

We use the following theorem-level inputs from earlier chapters:
\begin{enumerate}[label=\textup{(\roman*)}]
\item
      stabilized quadratic carrier
      \(\mathcal K\cong F^2/F^3\)
      (
      \cref{thm:interface-stabilized}
      );
\item
      one-dimensional canonical scalar channel on
      \(\mathcal K_{\mathbb R}\)
      (
      \cref{thm:quadratic-scalar,thm:spectral-scalar-identification}
      );
\item
      Hilbert realization branch of the stabilized quadratic carrier
      (
      \cref{ch:hilbert}
      );
\item
      phase realization branch of the same carrier
      (
      \cref{ch:em}
      );
\item
      macroscopic Einstein-branch scalar coefficient, read under the
      inherited second-jet faithfulness condition on the realized
      degree--\(2\) channel attached to the stabilized quadratic
      carrier
      (
      \cref{ch:einstein}
      );
\item
      canonical sector defect classes and scalar invariants
      \(\kappa_\sigma\in\mathcal K\),
      \(m_0(\sigma):=Q(\kappa_\sigma)\)
      (
      \cref{def:spectral-defect-class,def:spectral-scalar,thm:spectral-depth-separation}
      ).
\end{enumerate}

The sections below keep these inputs separated by role: carrier typing,
branch realization, and observer-side factorization.

% ============================================================
\section{Canonical sector scalar}
\label{sec:matter-realization-scalar}
% ============================================================

\begin{definition}[Canonical sector scalar]
\label{def:canonical-sector-scalar}
For a primitive excitation sector \(\sigma\), define
\[
  m_0(\sigma):=Q(\kappa_\sigma),
\]
where \(\kappa_\sigma\in\mathcal K\) is the canonical quadratic defect
class and \(Q\) is the unique admissible scalar channel on
\(\mathcal K_{\mathbb R}\).
\end{definition}

\begin{remark}[Chapter~19 role]
By \cref{thm:spectral-scalar-positive,thm:spectral-scalar-invariant,thm:spectral-depth-separation},
\(m_0(\sigma)\) is strictly positive, transport-invariant, and part of the
intrinsic sector-separation package.
\end{remark}

% ============================================================
\section{Typing lemmas for realization roles}
\label{sec:matter-realization-typing}
% ============================================================

\begin{lemma}[Carrier factorization of the canonical sector scalar]
\label{lem:matter-realization-factorization}
For every primitive excitation sector \(\sigma\), the scalar
\(m_0(\sigma)\) factors canonically as
\[
  \sigma
  \longmapsto
  \kappa_\sigma\in\mathcal K
  \hookrightarrow
  \mathcal K_{\mathbb R}
  \xrightarrow{\,Q\,}
  \mathbb R_{\ge0}.
\]
In particular, \(m_0\) is intrinsic and transport-invariant before any
observer-side interpretation map is introduced.
\end{lemma}

\begin{proof}
By \cref{def:spectral-defect-class}, each primitive sector has a canonical
defect class \(\kappa_\sigma\in\mathcal K\).
By \cref{thm:quadratic-scalar}, the scalar channel is a canonical positive
quadratic map
\(Q:\mathcal K_{\mathbb R}\to\mathbb R_{\ge0}\), and
\cref{def:spectral-scalar} defines
\(m_0(\sigma)=Q(\kappa_\sigma)\).
Transport invariance is exactly
\cref{thm:spectral-scalar-invariant}, and strict positivity for primitive
nontrivial sectors is \cref{thm:spectral-scalar-positive}.
Hence the displayed factorization and intrinsic status are fixed prior to
any external semantics.
\end{proof}

\begin{proposition}[Microscopic/macroscopic role split on one scalar channel]
\label{prop:matter-realization-role-split}
The canonical scalar channel admits two theorem-level realization roles:
\begin{enumerate}[label=\textup{(\roman*)}]
\item
      microscopic excitation-sector role, where
      \(m_0(\sigma)=Q(\kappa_\sigma)\) participates in primitive-sector
      classification;
\item
      macroscopic geometric role, where, under the inherited second-jet faithfulness condition, the same scalar channel is read as the
      Einstein-branch scalar coefficient on the realized degree--\(2\)
      channel attached to the stabilized quadratic carrier.
\end{enumerate}
These roles are type-distinct and are not identified with one another
without an additional realization map.
\end{proposition}

\begin{proof}
For \textup{(i)}, Chapter~19 defines \(m_0\) on primitive sectors and uses
it inside sector-separation structure
(\cref{def:spectral-scalar,thm:spectral-depth-separation}).
For \textup{(ii)}, Chapter~16 shows that, under the inherited
second-jet faithfulness condition, the macroscopic scalar coefficient is
read from the same canonical degree--\(2\) scalar channel on the
realized degree--\(2\) channel attached to the stabilized quadratic
carrier, up to normalization
(\cref{thm:einstein-beta-from-carrier,thm:einstein-classification}).

The codomains and logical roles differ: \textup{(i)} is a scalar invariant
attached to primitive excitation sectors, while \textup{(ii)} is a
coefficient in the macroscopic geometric operator classification.
Therefore no theorem-level identification between the two roles follows
without extra data.
\end{proof}

% ============================================================
\section{Branch selection}
\label{sec:matter-realization-branch}
% ============================================================

\begin{theorem}[Matter-sector branch selection]
\label{thm:matter-branch-selection}
Let \(\sigma\) be a primitive excitation sector and
\(m_0(\sigma)=Q(\kappa_\sigma)\) its canonical scalar invariant.
Then the admissible first realization role of \(m_0(\sigma)\) is through
the Hilbert--phase branch of the stabilized quadratic carrier.
Equivalently:
\begin{enumerate}[label=\textup{(\roman*)}]
\item
      as a sector observable, \(m_0(\sigma)\) is carried by the
      microscopic quadratic Hilbert law and its phase descent;
\item
      \(m_0(\sigma)\) is not, by default, a matter-sector observable in
      the macroscopic Einstein branch, where, under the inherited
      second-jet faithfulness condition, the same scalar channel is read
      as the Einstein-branch scalar coefficient on the realized
      degree--\(2\) channel attached to the stabilized quadratic
      carrier.
\end{enumerate}
\end{theorem}

\begin{proof}
By \cref{lem:matter-realization-factorization},
\(m_0(\sigma)\) is attached to the stabilized quadratic defect class
\(\kappa_\sigma\in\mathcal K\).
The first realization branch of this carrier is the microscopic
Hilbert-sector realization (\cref{thm:quadratic-scalar}) together with its
phase descent (\cref{thm:em-phase-reduction,thm:em-closedness-phase}).
So the intrinsic excitation-sector scalar is already realized on the
Hilbert--phase branch.

By \cref{prop:matter-realization-role-split}, under the inherited
second-jet faithfulness condition the Einstein branch reads the same
scalar channel as macroscopic geometric coefficient data on the
realized degree--\(2\) channel attached to the stabilized quadratic
carrier, rather than as primitive-sector observable data.
Since Chapter~19 employs
\(m_0(\sigma)\) as part of primitive excitation-sector organization
(\cref{thm:spectral-depth-separation}), its first admissible realization
role is the microscopic Hilbert--phase role.
Hence \textup{(i)} and \textup{(ii)} follow.
\end{proof}

% ============================================================
\section{Sector-observable priority}
\label{sec:matter-realization-priority}
% ============================================================

\begin{theorem}[Sector-observable priority]
\label{thm:sector-observable-priority}
\label{thm:sector-observable}
For every primitive excitation sector \(\sigma\),
\[
  m_0(\sigma)=Q(\kappa_\sigma)
\]
is, prior to any observer-side semantics, an intrinsic sector observable.
\end{theorem}

\begin{proof}
By \cref{lem:matter-realization-factorization}, the map
\(\sigma\mapsto m_0(\sigma)\) is defined purely from
\(\kappa_\sigma\in\mathcal K\) and the intrinsic scalar channel \(Q\), and
is transport-invariant before any observer semantics.
Chapter~19 then uses this intrinsic scalar inside sector-organization data
(\cref{thm:spectral-depth-separation}).
By \cref{thm:matter-branch-selection}, this scalar is first realized on the
Hilbert--phase branch at sector level.
Therefore its status as a sector observable is fixed internally, prior to
any observer-side interpretation.
\end{proof}

% ============================================================
\section{Realization boundary}
\label{sec:matter-realization-boundary}
% ============================================================

\begin{corollary}[Mass-realization boundary]
\label{cor:mass-boundary}
\label{cor:boundary}
Any identification of the intrinsic sector scalar \(m_0(\sigma)\), or of
the Chapter~20 closed-system mass coordinate, with an observer-side measured
mass observable requires a separate semantic realization theorem.
It is not part of the intrinsic theorem package of
\cref{ch:causality,ch:hilbert,ch:einstein,ch:connection,ch:em,ch:matter,ch:numbers-shell-closure}.
\end{corollary}

\begin{proof}
By \cref{thm:sector-observable-priority},
\(m_0(\sigma)\) is theorem-level intrinsic sector data.
By \cref{thm:matter-branch-selection} and
\cref{prop:matter-realization-role-split}, this intrinsic datum has a
canonical microscopic role and is not automatically identified with
the Einstein-branch scalar coefficient read, under the inherited
second-jet faithfulness condition, on the realized degree--\(2\)
channel attached to the stabilized quadratic carrier.

Hence any external measured-mass interpretation must introduce an additional
semantic readout map, schematically
\[
  I_{\mathrm{mass}}:
  m_0(\sigma)
  \longmapsto
  m_{\mathrm{obs}}(\sigma),
\]
not determined by the intrinsic package listed in the statement.
Chapter~20 proves the internal quartic mass-coordinate bridge
(\cref{thm:numbers-mass-calibration,rem:numbers-forced-vs-calibrated}); what is
separate is the observer-side identification of that internal coordinate with
a measured mass convention.  Therefore external measured-mass interpretation
remains a semantic realization theorem, not an internal axiom.
\end{proof}

% ============================================================
\section{Operator-first spectral realization gap}
\label{sec:matter-realization-gap}
% ============================================================

The next theorem-level strengthening is not orthogonality as a starting
input.
The exact remaining direction is to construct the Chapter~19 sector scalar
invariant as a self-adjoint operator on the Hilbert realization branch.
Orthogonality is then a consequence of spectral structure.

From the existing chapter stack, we have the following proved data:
\begin{itemize}
\item stabilized quadratic carrier
      \(\mathcal K\cong F^2/F^3\)
      (\cref{thm:interface-stabilized});
\item unique one-dimensional scalar channel on
      \(\mathcal K_{\mathbb R}\)
      (\cref{thm:quadratic-scalar,thm:spectral-scalar-identification});
\item Hilbert realization of the quadratic carrier
      (\cref{ch:hilbert});
\item intrinsic excitation-sector data
      \(\kappa_\sigma\),
      \(m_0(\sigma)=Q(\kappa_\sigma)\),
      \(k_\sigma\), which separates primitive sectors inside a fixed gauge
      class (\cref{thm:spectral-depth-separation}).
\end{itemize}

None of these, by itself, yields a canonical densely defined positive
self-adjoint operator on the Hilbert realization whose sector eigenvalues
are the intrinsic scalars \(m_0(\sigma)\).

\begin{proposition}[Nature of the realization gap]
\label{prop:matter-realization-gap}
Intrinsic sector separation does not by itself produce the Hilbert-branch
sector scalar operator.
\end{proposition}

\begin{proof}
Sector separation is proved at classification level in the stabilized
quadratic carrier, through invariants such as
\(\kappa_\sigma\), \(m_0(\sigma)\), and \(k_\sigma\)
(\cref{def:spectral-defect-class,def:spectral-scalar,thm:spectral-depth-separation}).
Operator realization is a statement about unbounded self-adjoint structure on
the Hilbert branch.
No theorem in the current chain yet constructs the canonical positive
self-adjoint operator whose sector eigenvalues are the intrinsic scalars.
Hence the remaining gap is an operator-construction gap.
\end{proof}

\begin{theorem}[Target operator realization of the sector scalar invariant]
\label{thm:hilbert-realization-sector-scalar}
\label{thm:sector-observable-realization}
Let \(H_{\mathrm{quad}}\) be the Hilbert realization of the stabilized
quadratic carrier \(\mathcal K\simeq F^2/F^3\) from Chapter~15.
The exact remaining realization-level operator target is to prove that
there exists a canonical densely defined positive self-adjoint operator
\[
  \mathsf S : D(\mathsf S)\subset H_{\mathrm{quad}} \to H_{\mathrm{quad}}
\]
such that for every primitive excitation sector \(\sigma\) one has a nonzero
vector \(\psi_\sigma \in D(\mathsf S)\) satisfying
\[
  \mathsf S \psi_\sigma = m_0(\sigma)\psi_\sigma,
  \qquad
  m_0(\sigma)=Q(\kappa_\sigma).
\]
\end{theorem}

\begin{remark}[Status of the target theorem]
\label{rem:matter-realization-sector-observable-status}
\Cref{thm:hilbert-realization-sector-scalar} records the exact remaining
realization-level operator statement to be discharged.
It is not part of the current intrinsic theorem package proved so far.
\end{remark}

\begin{lemma}[Target form-typing input for the sector operator route]
\label{lem:sector-quadratic-form}
A sufficient intermediate target for
\cref{thm:hilbert-realization-sector-scalar}
is to prove that there exists a densely defined positive sesquilinear form
\[
  \mathfrak q : D(\mathfrak q)\times D(\mathfrak q)\to \mathbb C
\]
on \(H_{\mathrm{quad}}\) such that for every primitive excitation sector
\(\sigma\),
\[
  \mathfrak q(\psi_\sigma,\psi_\sigma)=m_0(\sigma)\|\psi_\sigma\|^2.
\]
\end{lemma}

\begin{remark}[Conditional candidate for the form route]
\label{rem:matter-realization-form-candidate}
If the later sector-to-Hilbert embedding and form-typing targets are
discharged, then the natural candidate for the form in
\cref{lem:sector-quadratic-form}
is the defect-carrier expression
\[
  \mathfrak q(\psi,\phi):=\langle\kappa(\psi),\kappa(\phi)\rangle_Q,
\]
where \(\kappa(\psi)\) denotes the quadratic defect representative
attached to the embedded sector vector \(\psi\), and
\(\langle\cdot,\cdot\rangle_Q\)
is the polarized form induced by the scalar channel \(Q\)
(\cref{prop:hilbert-real-inner}) together with its later
sesquilinear Hilbert-branch extension.
This remark records only that candidate route;
it does not assert that the required embedding, form-typing, or
closability input has already been proved.
\end{remark}

\begin{proposition}[Closed-form route to the sector operator]
\label{prop:matter-realization-form-route}
Assume \cref{lem:sector-quadratic-form} and assume the form \(\mathfrak q\) is closed
or closable.
Then the representation theorem for closed positive forms yields a canonical
positive self-adjoint operator
\[
  \mathsf S:D(\mathsf S)\subset H_{\mathrm{quad}}\to H_{\mathrm{quad}}
\]
satisfying
\[
  \mathfrak q(\psi,\phi)=\langle\psi,\mathsf S\phi\rangle
\]
for all \(\psi\in D(\mathfrak q)\), \(\phi\in D(\mathsf S)\).
\end{proposition}

\begin{proof}
If \(\mathfrak q\) is closed and positive, Kato's representation theorem for
closed positive forms assigns
a unique positive self-adjoint operator \(\mathsf S\) whose form domain is
\(D(\mathfrak q)\) and whose form identity is exactly the displayed equality.
If \(\mathfrak q\) is closable, apply the same theorem to its closure
\cite{Kato1995}.
Canonicity follows from uniqueness in the representation theorem.
\end{proof}

% ============================================================
\subsection{Proof checklist for the forced-upgrade roadmap}
\label{subsec:matter-realization-roadmap-checklist}
% ============================================================

This subsection maps the roadmap stubs of
\cref{subsec:shell-forced-upgrade-roadmap}
to exact proof slots inside the Chapter~22 operator route.

\begin{proposition}[Later roadmap-target bookkeeping for the Chapter~22 operator route]
\label{prop:matter-roadmap-stub-slot-map}
If the later roadmap targets recorded in
\cref{subsec:shell-forced-upgrade-roadmap}
were discharged, then they would feed this chapter as follows:
\begin{enumerate}[label=\textup{(\roman*)}]
\item
      \cref{thm:roadmap-sector-to-hilbert-embedding}
      would supply the canonical sector-to-Hilbert embedding input needed
      to
      place sector vectors \(\psi_\sigma\in H_{\mathrm{quad}}\) in the
      operator-construction route of \cref{sec:matter-realization-gap};
\item
      \cref{lem:roadmap-embedded-sector-form}
      would supply the form-typing input slot formalized by
      \cref{lem:sector-quadratic-form};
\item
      \cref{thm:roadmap-closed-form-closability}
      would discharge the closed/closable hypothesis required by
      \cref{prop:matter-realization-form-route}.
\end{enumerate}
\end{proposition}

\begin{proof}
Immediate from the repaired target status of
\cref{lem:sector-quadratic-form,prop:matter-realization-form-route}
and the later-upgrade bookkeeping already recorded in
\cref{subsec:shell-forced-upgrade-roadmap,prop:roadmap-implies-ch22-operator-target}.
\end{proof}

\begin{corollary}[Later checklist discharge promotes the operator target]
\label{cor:matter-roadmap-checklist-to-operator-target}
If the three roadmap targets named in
\cref{prop:matter-roadmap-stub-slot-map}
are later proved from existing theorem-level inputs without new
observer-side assumptions, then that later discharge promotes the
Chapter~22 operator target
\cref{thm:hilbert-realization-sector-scalar}
through the closed-form route.
\end{corollary}

\begin{proof}
Combine
\cref{prop:matter-roadmap-stub-slot-map,prop:matter-realization-form-route,prop:roadmap-implies-ch22-operator-target}.
\end{proof}

\begin{lemma}[Sector orthogonality from operator realization]
\label{lem:matter-realization-sector-orthogonality}
Let \(\mathsf S\) be self-adjoint on \(H_{\mathrm{quad}}\), and suppose
\[
  \mathsf S\psi_\sigma = m_0(\sigma)\psi_\sigma,
  \qquad
  \mathsf S\psi_\tau = m_0(\tau)\psi_\tau.
\]
Then
\[
  m_0(\sigma)\neq m_0(\tau)
  \Longrightarrow
  \langle\psi_\sigma,\psi_\tau\rangle=0.
\]
\end{lemma}

\begin{proof}
Self-adjointness gives
\[
  \langle \mathsf S\psi_\sigma,\psi_\tau\rangle
  =
  \langle \psi_\sigma,\mathsf S\psi_\tau\rangle.
\]
Hence
\[
  m_0(\sigma)\langle\psi_\sigma,\psi_\tau\rangle
  =
  m_0(\tau)\langle\psi_\sigma,\psi_\tau\rangle,
\]
so
\[
  (m_0(\sigma)-m_0(\tau))\langle\psi_\sigma,\psi_\tau\rangle=0.
\]
If \(m_0(\sigma)\neq m_0(\tau)\), then
\(\langle\psi_\sigma,\psi_\tau\rangle=0\).
\end{proof}

\begin{corollary}[Consequences conditional on sector-observable realization]
\label{cor:matter-realization-sector-observable-consequences}
Assume \cref{thm:hilbert-realization-sector-scalar} and
\[
  H_{\mathrm{quad}}=
  \overline{\operatorname{span}\{\psi_\sigma:
  \sigma\in\mathsf S_{\mathrm{prim}}\}}.
\]
Then:
\begin{enumerate}[label=\textup{(\roman*)}]
\item
      \(m_0(\sigma)\neq m_0(\tau)
      \Rightarrow \langle\psi_\sigma,\psi_\tau\rangle=0\);
\item
      the sector eigenvectors generate a genuine spectral family for
      \(\mathsf S\);
\item
      canonical spectral projections for \(\mathsf S\) exist and determine
      the sector decomposition of the Hilbert branch.
\end{enumerate}
\end{corollary}

\begin{proof}
Item \textup{(i)} is \cref{lem:matter-realization-sector-orthogonality}.
Items \textup{(ii)} and \textup{(iii)} are standard consequences of the
spectral theorem for self-adjoint operators once the eigenvector family is
complete in the stated closure \cite{Kato1995}.
\end{proof}

In this way, orthogonality, spectral families, and projections appear as
consequences of operator realization, not as independent starting input.
The remaining realization step is therefore exact and isolated:
construct the sector operator by closed positive form methods.

% ============================================================
\section{Observer-map factorization algebra}
\label{sec:matter-realization-factorization-algebra}
% ============================================================

Write
\[
  \mathsf S_{\mathrm{prim}}
  :=
  \{\sigma:\sigma\text{ primitive excitation sector}\},
  \qquad
  m_0:\mathsf S_{\mathrm{prim}}\to\mathbb R_{\ge0},
  \ \sigma\mapsto Q(\kappa_\sigma).
\]

The goal of this section is formal and minimal: isolate exactly when an
observer-side assignment descends through the intrinsic scalar channel.

\begin{definition}[Observer-side mass assignment]
\label{def:matter-realization-observer-map}
An observer-side mass assignment is any map
\[
  M_{\mathrm{obs}}:\mathsf S_{\mathrm{prim}}\to\mathbb R_{>0}.
\]
\end{definition}

\begin{theorem}[Factorization criterion through the intrinsic scalar]
\label{thm:matter-realization-factorization-criterion}
Let \(M_{\mathrm{obs}}\) be an observer-side mass assignment.
The following are equivalent:
\begin{enumerate}[label=\textup{(\roman*)}]
\item there exists a map
      \(\Lambda:\mathrm{Im}(m_0)\to\mathbb R_{>0}\)
      such that
      \[
        M_{\mathrm{obs}}=\Lambda\circ m_0;
      \]
\item for all
      \(\sigma_1,\sigma_2\in\mathsf S_{\mathrm{prim}}\),
      \[
        m_0(\sigma_1)=m_0(\sigma_2)
        \Longrightarrow
        M_{\mathrm{obs}}(\sigma_1)=M_{\mathrm{obs}}(\sigma_2).
      \]
\end{enumerate}
When these conditions hold, \(\Lambda\) is unique.
\end{theorem}

\begin{proof}
\textup{(i)}\(\Rightarrow\)\textup{(ii)}:
if \(M_{\mathrm{obs}}=\Lambda\circ m_0\) and
\(m_0(\sigma_1)=m_0(\sigma_2)\), then
\[
  M_{\mathrm{obs}}(\sigma_1)
  =\Lambda(m_0(\sigma_1))
  =\Lambda(m_0(\sigma_2))
  =M_{\mathrm{obs}}(\sigma_2).
\]

\textup{(ii)}\(\Rightarrow\)\textup{(i)}:
define
\(\Lambda:\mathrm{Im}(m_0)\to\mathbb R_{>0}\)
by
\[
  \Lambda(x):=M_{\mathrm{obs}}(\sigma)
  \quad
  \text{for any }\sigma\text{ with }m_0(\sigma)=x.
\]
Condition \textup{(ii)} makes this well-defined.
Then for every \(\sigma\),
\[
  (\Lambda\circ m_0)(\sigma)=\Lambda(m_0(\sigma))=M_{\mathrm{obs}}(\sigma),
\]
so \(M_{\mathrm{obs}}=\Lambda\circ m_0\).
For uniqueness, if \(\Lambda_1\circ m_0=\Lambda_2\circ m_0\), then for
every \(x\in\mathrm{Im}(m_0)\) choose \(\sigma\) with
\(m_0(\sigma)=x\), giving
\(\Lambda_1(x)=M_{\mathrm{obs}}(\sigma)=\Lambda_2(x)\).
Hence \(\Lambda_1=\Lambda_2\).
\end{proof}

\begin{corollary}[Non-injectivity obstruction]
\label{cor:matter-realization-noninjective-obstruction}
Assume there exist
\(\sigma_1\neq\sigma_2\) with
\[
  m_0(\sigma_1)=m_0(\sigma_2),
  \qquad
  M_{\mathrm{obs}}(\sigma_1)\neq M_{\mathrm{obs}}(\sigma_2).
\]
Then \(M_{\mathrm{obs}}\) cannot factor through \(m_0\).
Hence any such assignment requires additional sector observables beyond
the intrinsic scalar channel.
\end{corollary}

\begin{proof}
If \(M_{\mathrm{obs}}\) factored through \(m_0\),
\cref{thm:matter-realization-factorization-criterion} would force equality
\(M_{\mathrm{obs}}(\sigma_1)=M_{\mathrm{obs}}(\sigma_2)\), contradicting
the hypothesis.
The final statement is immediate.
\end{proof}

\begin{proposition}[Two-point affine calibration on \(\mathrm{Im}(m_0)\)]
\label{prop:matter-realization-affine-two-point}
Fix two distinct intrinsic scalar values
\(x_1,x_2\in\mathrm{Im}(m_0)\) and target values
\(y_1,y_2\in\mathbb R\).
There is a unique affine map
\(
  \Lambda_{\mathrm{aff}}(x)=ax+b
\)
satisfying
\[
  \Lambda_{\mathrm{aff}}(x_1)=y_1,
  \qquad
  \Lambda_{\mathrm{aff}}(x_2)=y_2,
\]
namely
\[
  a=\frac{y_2-y_1}{x_2-x_1},
  \qquad
  b=\frac{x_2y_1-x_1y_2}{x_2-x_1}.
\]
\end{proposition}

\begin{proof}
The constraints are the linear system
\[
  ax_1+b=y_1,
  \qquad
  ax_2+b=y_2.
\]
Subtract to obtain
\(
  a(x_2-x_1)=y_2-y_1
\), hence the stated formula for \(a\).
Substituting into either equation gives the stated formula for \(b\).
Because \(x_2\neq x_1\), the system has unique solution.
\end{proof}

\begin{remark}[Boundary-data interpretation]
\label{rem:matter-realization-factorization-boundary}
\Cref{thm:matter-realization-factorization-criterion} and
\cref{prop:matter-realization-affine-two-point} formalize the same logical
boundary recorded in \cref{cor:mass-boundary}:
the intrinsic sector scalar defines the factorization channel, while the
choice of observer map \(\Lambda\) (or calibration parameters) is additional
realization data.
\end{remark}

% ============================================================
\section{Conclusion}
\label{sec:matter-realization-conclusion}
% ============================================================

The canonical scalar invariant of Chapter~19 is now placed at its proper
realization level.
It is:
\begin{itemize}
\item intrinsic;
\item sector-theoretic;
\item first realized through the Hilbert--phase branch;
\item constrained by an explicit observer-map factorization criterion;
\item not yet an observer-side mass observable without extra calibration.
\end{itemize}

The remaining intrinsic realization step is explicit and isolated:
construct the Hilbert-branch sector operator by the closed positive
form route recorded in
\cref{thm:hilbert-realization-sector-scalar,lem:sector-quadratic-form,prop:matter-realization-form-route}.
Any later physical-mass identification would still require a separate
observer-side realization theorem or calibration input beyond that
operator construction.

The following chapter uses the completed realization stack as input for a
different question. It does not alter the scalar-sector boundary proved
here. Instead, it asks how much of the state-side limit observable
algebra is readable by a smooth realization, and which derivations of
the readable sector correspond to intrinsic automorphisms of the
realized connection.
